\documentclass[11pt, a4paper]{article}

% --- PACOTES PRINCIPAIS ---
\usepackage[utf8]{inputenc}
\usepackage[T1]{fontenc}
\usepackage{geometry}
\geometry{a4paper, top=2.5cm, bottom=2.5cm, left=2.0cm, right=2.0cm}
\usepackage{amsmath, amssymb, amsfonts}
\usepackage{booktabs}
\usepackage{tabularx}
\usepackage{graphicx}
\usepackage{xcolor}
\usepackage{enumitem}
\usepackage{tcolorbox}
\usepackage{fancyhdr}
\usepackage{float}
\usepackage[colorlinks=true, linkcolor=ufopred, citecolor=deepnavy, urlcolor=ufopred]{hyperref}

\graphicspath{{figures/}}

% nomes em português sem depender do babel
\renewcommand{\figurename}{Figura}
\renewcommand{\tablename}{Tabela}
\renewcommand{\refname}{Referências}

% --- CORES INSTITUCIONAIS UFOP ---
\definecolor{ufopred}{RGB}{140, 29, 36}
\definecolor{deepnavy}{RGB}{20, 35, 60}
\definecolor{lightgray}{RGB}{245, 247, 250}
\definecolor{darkslate}{RGB}{30, 40, 50}

% --- CONFIGURAÇÃO DE CABEÇALHO E RODAPÉ ---
\pagestyle{fancy}
\fancyhf{}
\setlength{\headheight}{14pt}
\rhead{\textcolor{darkslate}{\small \textbf{BCC327 -- Computação Gráfica} | TP1}}
\lhead{\textcolor{darkslate}{\small Brenda Gabrielle Alves Nascimento | 24.1.4011}}
\cfoot{\thepage}

\tcbuselibrary{skins, breakable}

% figura lado a lado com legenda (a), (b)...
\newcommand{\painel}[3]{%
  \begin{minipage}[t]{#1\textwidth}\centering
    \includegraphics[width=\linewidth]{#2}\\[-0.1cm]
    {\small #3}
  \end{minipage}}

\begin{document}

% --- TÍTULO E AUTORES ---
\begin{center}
  {\LARGE \bfseries \textcolor{ufopred}{Visualização Geométrica, Rasterização e Recorte de Árvores Arteriais CCO}}\\[0.4cm]
  {\large \textbf{Brenda Gabrielle Alves Nascimento}$^{1}$ (24.1.4011)}\\[0.2cm]
  {\small $^{1}$Departamento de Computação (DECOM), Universidade Federal de Ouro Preto (UFOP)\\[0.1cm]
  BCC327 -- Computação Gráfica, 2026.2 | Prof. Dr. Rafael Alves Bonfim de Queiroz}\\[0.4cm]
\end{center}

% --- CAIXA DE RESUMO E PALAVRAS-CHAVE ---
\begin{tcolorbox}[colback=lightgray, colframe=deepnavy, arc=3pt, title=\textbf{RESUMO}, fonttitle=\bfseries]
\small
\textbf{Contexto:} Árvores arteriais sintéticas geradas pelo método \emph{Constrained Constructive Optimization} (CCO) são descritas por pontos, segmentos e o raio de cada vaso, e precisam ser convertidas em imagem para análise visual.
\textbf{Objetivo:} Desenvolver uma aplicação de Computação Gráfica que visualize os modelos CCO 2D e 3D e demonstre, com implementação própria, os algoritmos clássicos envolvidos na construção da imagem.
\textbf{Metodologia:} A aplicação, em Python com pygame-ce, PyOpenGL e NumPy, tem um visualizador com leitor VTK, mapas de cores, faixas e tubos proporcionais ao raio, matrizes próprias de transformação, câmera e projeção, e MSAA; e um modo didático cujo framebuffer é escrito apenas pelos algoritmos do grupo: Bresenham, DDA, Scanline Fill, supersampling, Cohen--Sutherland, Liang--Barsky, Sutherland--Hodgman e Weiler--Atherton.
\textbf{Resultados:} Os 12 modelos foram lidos corretamente; as matrizes coincidem com as funções equivalentes do OpenGL; os algoritmos foram validados contra testes de referência sem divergências; e as comparações visuais evidenciam projeções, aliasing e os limites do Sutherland--Hodgman com polígonos côncavos.
\textbf{Conclusão:} A separação entre visualizador e modo didático permitiu aliar uma visualização fluida dos modelos à demonstração transparente, pixel a pixel, de cada algoritmo.

\vspace{0.2cm}
\textbf{Palavras-Chave:} Computação Gráfica; Método CCO; Rasterização; Recorte; Antialiasing.
\end{tcolorbox}

\vspace{0.3cm}

% ======================================================================
\section{Introdução}

O método CCO, introduzido por Schreiner e Buxbaum~\cite{schreiner1993}, gera árvores vasculares sintéticas que suprem um domínio de perfusão sob restrições fisiológicas. Cada modelo é um grafo em árvore: pontos no espaço, segmentos que ligam pares de pontos e um raio associado a cada segmento. Para que esses modelos sejam inspecionados, é preciso transformá-los em imagem, o que envolve todas as etapas clássicas do \emph{pipeline} gráfico.

Este trabalho desenvolve uma aplicação para visualizar, transformar, rasterizar e recortar os modelos CCO 2D e 3D fornecidos em formato VTK. O foco não é apenas usar uma biblioteca gráfica, mas implementar e demonstrar os algoritmos clássicos. As contribuições são:
\begin{itemize}[leftmargin=*]
  \item um leitor VTK próprio e um visualizador com linhas, faixas 2D e tubos 3D codificados por cor e espessura em função do raio;
  \item matrizes próprias de translação, rotação, escala, câmera e projeções ortográfica e perspectiva, validadas contra o OpenGL;
  \item um modo didático com framebuffer próprio para Bresenham, DDA, Scanline Fill e supersampling;
  \item implementação e comparação de Cohen--Sutherland, Liang--Barsky, Sutherland--Hodgman e, como diferencial, Weiler--Atherton.
\end{itemize}

% ======================================================================
\section{Fundamentação Teórica}

\subsection{Modelos CCO e formato VTK}
Os modelos estão no formato VTK legado ASCII, com \texttt{DATASET POLYDATA}. A seção \texttt{POINTS} lista as coordenadas $(x, y, z)$; a seção \texttt{LINES} lista as células, cada uma no formato \texttt{2 i j}, ligando o ponto $i$ ao ponto $j$; e a seção \texttt{CELL\_DATA} traz o escalar \texttt{raio}, um valor por segmento. Como o raio é atributo de célula, toda a cor e a espessura de um vaso derivam de um único valor. As coordenadas estão em metros (domínio com cerca de 5\,cm de raio) e o raio em milímetros.

\subsection{Modelo de cor}
Um mapa contínuo de cores associa a cada escalar normalizado $t = (r - r_{\min})/(r_{\max} - r_{\min}) \in [0, 1]$ uma cor RGB obtida por interpolação linear entre cores de controle. Foram usados dois mapas. O \emph{jet} (arco-íris) reproduz o estilo das figuras de referência e tem alto contraste, mas não é perceptualmente uniforme: o ciano e o amarelo criam fronteiras aparentes inexistentes nos dados~\cite{borland2007}. O \emph{viridis} é perceptualmente uniforme e legível por daltônicos e em escala de cinza. O modelo RGB foi adotado por ser o espaço nativo do framebuffer.

\subsection{Transformações geométricas e câmera}
Em coordenadas homogêneas, um ponto $\mathbf{p} = (x, y, z, 1)^T$ é transformado por matrizes $4 \times 4$, o que permite expressar translação como produto matricial:
\begin{equation}
  T = \begin{bmatrix} 1&0&0&t_x\\ 0&1&0&t_y\\ 0&0&1&t_z\\ 0&0&0&1 \end{bmatrix},\quad
  S = \begin{bmatrix} s_x&0&0&0\\ 0&s_y&0&0\\ 0&0&s_z&0\\ 0&0&0&1 \end{bmatrix},\quad
  R_z(\theta) = \begin{bmatrix} \cos\theta&-\sin\theta&0&0\\ \sin\theta&\cos\theta&0&0\\ 0&0&1&0\\ 0&0&0&1 \end{bmatrix}.
\end{equation}
O vértice percorre $\mathbf{p}_{\text{recorte}} = P \cdot V \cdot M \cdot \mathbf{p}_{\text{modelo}}$, em que $M$ leva o modelo ao \textbf{espaço do mundo}, $V$ leva o mundo ao \textbf{sistema de coordenadas da câmera} e $P$ aplica a projeção. O espaço do mundo é o referencial fixo da cena, comum a todos os objetos. O sistema da câmera é a mesma cena descrita do ponto de vista do observador: a câmera fica na origem, olhando para $-z$. A matriz $V$ é montada a partir do olho $\mathbf{e}$, do alvo $\mathbf{a}$ e de um vetor ``cima'' $\mathbf{c}$, que definem a base ortonormal da câmera
\begin{equation}
  \mathbf{w} = \frac{\mathbf{e} - \mathbf{a}}{\|\mathbf{e} - \mathbf{a}\|},\qquad
  \mathbf{u} = \frac{\mathbf{c} \times \mathbf{w}}{\|\mathbf{c} \times \mathbf{w}\|},\qquad
  \mathbf{v} = \mathbf{w} \times \mathbf{u},
\end{equation}
e $V$ translada o olho para a origem e reescreve cada ponto nessa base.

Na \textbf{projeção ortográfica} os raios de projeção são paralelos: a matriz apenas mapeia uma caixa no cubo $[-1,1]^3$, e o tamanho aparente não depende da distância. Na \textbf{projeção perspectiva}, a última linha da matriz é $(0, 0, -1, 0)$, que copia a profundidade para $w$; na divisão por $w$, objetos distantes ficam menores e retas paralelas convergem.

\subsection{Rasterização de segmentos}
O algoritmo de Bresenham~\cite{bresenham1965} rasteriza um segmento usando apenas aritmética inteira. No caso $0 \le m \le 1$, $x$ avança sempre uma unidade e a variável de decisão $p_k$, proporcional à diferença entre as distâncias da reta ideal aos dois pixels candidatos, escolhe entre manter ou incrementar $y$:
\begin{equation}
  p_0 = 2\Delta y - \Delta x,\qquad
  p_{k+1} = \begin{cases} p_k + 2\Delta y, & p_k < 0 \\ p_k + 2\Delta y - 2\Delta x \ \ (y \leftarrow y+1), & p_k \ge 0. \end{cases}
\end{equation}
Os demais octantes são obtidos trocando os papéis de $x$ e $y$ e os sinais dos incrementos. O DDA, implementado para comparação, soma incrementos reais $\Delta y / \Delta x$ e arredonda cada posição, exigindo ponto flutuante a cada passo.

\subsection{Preenchimento de polígonos por Scanline}
O Scanline Fill~\cite{foley1990} percorre o polígono linha a linha. Uma \emph{tabela de arestas} (ET) guarda cada aresta não horizontal no balde da linha onde começa, com $(y_{\max}, x, 1/m)$. Em cada linha, as arestas que começam entram na \emph{tabela de arestas ativas} (AET), as que terminaram saem, os $x$ são ordenados e preenchidos aos pares (regra par-ímpar), e cada $x$ avança $1/m$ para a linha seguinte, aproveitando a coerência entre linhas. O intervalo semiaberto $[y_{\min}, y_{\max})$ faz com que um vértice compartilhado seja contado uma única vez.

\subsection{Aliasing e antialiasing}
O aliasing surge porque um pixel só pode estar aceso ou apagado: bordas inclinadas viram degraus e detalhes menores que um pixel geram padrões espúrios. No \textbf{supersampling} (SSAA), a cena é renderizada numa grade $k$ vezes maior em cada direção e reduzida pela média de cada bloco $k \times k$, de modo que a cor final é proporcional à cobertura. No \textbf{MSAA}, cada pixel armazena várias amostras de cobertura e profundidade, mas o fragmento é sombreado uma única vez; na resolução, as amostras são promediadas. O custo é muito menor que o do SSAA e o efeito se concentra nas bordas dos objetos.

\subsection{Recorte de segmentos}
O \textbf{Cohen--Sutherland}~\cite{newman1979} atribui a cada extremidade um código de 4 bits que indica em quais semiplanos externos à janela o ponto está. Se $c_0 \lor c_1 = 0$, o segmento é aceito inteiro; se $c_0 \land c_1 \neq 0$, as duas extremidades estão do mesmo lado de fora e o segmento é rejeitado; caso contrário, uma extremidade externa é substituída pela interseção com uma borda e o teste se repete. O \textbf{Liang--Barsky}~\cite{liang1984} usa a forma paramétrica $P(t) = P_0 + t(P_1 - P_0)$, $t \in [0,1]$, e as desigualdades $p_k t \le q_k$, com $p = (-\Delta x, \Delta x, -\Delta y, \Delta y)$ e $q = (x_0 - x_{\min}, x_{\max} - x_0, y_0 - y_{\min}, y_{\max} - y_0)$. Para $p_k < 0$ a reta entra pela borda e $t_0 = \max(t_0, q_k/p_k)$; para $p_k > 0$ ela sai e $t_1 = \min(t_1, q_k/p_k)$; se $t_0 > t_1$, nada resta. As interseções só são calculadas no final.

O \textbf{Cyrus--Beck}~\cite{cyrus1978}, apenas discutido, generaliza a ideia paramétrica para \emph{qualquer janela convexa}: para cada aresta, com normal $\mathbf{n}_i$ e ponto $\mathbf{f}_i$, resolve $\mathbf{n}_i \cdot (P(t) - \mathbf{f}_i) = 0$ e classifica a interseção como entrada ou saída pelo sinal de $\mathbf{n}_i \cdot (P_1 - P_0)$. Sua principal diferença em relação aos dois anteriores é não exigir janela retangular alinhada aos eixos; o Liang--Barsky é justamente o Cyrus--Beck especializado para esse caso, em que os produtos escalares se reduzem aos $p_k$ e $q_k$.

\subsection{Recorte de polígonos}
O \textbf{Sutherland--Hodgman}~\cite{sutherland1974} trata a janela convexa como interseção de semiplanos e recorta o polígono contra cada aresta em sequência. Para cada aresta $S \to P$ do polígono: se ambos estão dentro, emite $P$; se só $S$ está dentro, emite a interseção $I$; se ambos estão fora, nada; se só $P$ está dentro, emite $I$ e $P$. Ele sempre produz \emph{um único} polígono; quando o resultado correto tem várias componentes, elas ficam ligadas por arestas degeneradas sobre a borda da janela. O \textbf{Weiler--Atherton}~\cite{weiler1977} insere as interseções nas listas de vértices do sujeito e da janela, classifica-as como entrada ou saída e percorre as duas fronteiras alternadamente, produzindo uma lista de polígonos. Ele é mais adequado quando o sujeito ou a janela são côncavos, quando o resultado tem múltiplas regiões ou quando há interseções complexas entre as fronteiras.

% ======================================================================
\section{Metodologia e Arquitetura da Aplicação}

A aplicação foi escrita em Python 3 com pygame-ce 2.5.2 (janela e entrada), PyOpenGL 3.1.7 e NumPy 2.1.3. Ela é composta por dois programas que compartilham os módulos da Tabela~\ref{tab:modulos}: o \textbf{visualizador}, que usa OpenGL para desenhar os modelos com as matrizes e a geometria calculadas pelo grupo, e o \textbf{modo didático}, uma matriz NumPy de $80 \times 60$ pixels, ampliada na tela, na qual apenas os algoritmos do grupo escrevem. O pygame-ce somente amplia essa matriz, de modo que nenhum pixel do modo didático é produzido pela biblioteca.

\begin{table}[H]
  \centering
  \small
  \caption{Módulos da aplicação e origem de cada funcionalidade.}
  \label{tab:modulos}
  \begin{tabularx}{\textwidth}{l X l}
    \toprule
    \textbf{Módulo} & \textbf{Conteúdo} & \textbf{Origem} \\
    \midrule
    \texttt{leitor\_vtk.py} & Leitor de VTK legado (POINTS, LINES, CELL\_DATA) & grupo \\
    \texttt{cores.py} & Mapas jet e viridis por interpolação linear & grupo \\
    \texttt{geometria.py} & Faixas 2D (retângulos e discos) e tubos 3D (cilindros, esferas e normais) & grupo \\
    \texttt{transformacoes.py} & Translação, rotação, escala, câmera, projeções & grupo \\
    \texttt{rasterizacao.py} & Bresenham, DDA e Scanline Fill & grupo \\
    \texttt{antialiasing.py} & Supersampling com filtro caixa & grupo \\
    \texttt{recorte.py} & Cohen--Sutherland, Liang--Barsky, Sutherland--Hodgman, Weiler--Atherton & grupo \\
    \texttt{visualizador.py} & Rasterização dos triângulos, profundidade, iluminação, MSAA, VBOs & OpenGL \\
    \texttt{modo\_didatico.py} & Framebuffer próprio e interações; janela e texto & grupo / pygame-ce \\
    \bottomrule
  \end{tabularx}
\end{table}

\textbf{Primitivas.} Além de linhas, cada segmento 2D vira uma faixa de largura $2r$ (dois triângulos), com um disco em cada extremidade que arredonda as junções. Em 3D, cada segmento vira um cilindro de 12 lados construído a partir da base ortonormal $(\mathbf{d}, \mathbf{e}_1, \mathbf{e}_2)$ do segmento, com esferas nos pontos da árvore; cada esfera recebe o raio e a cor do maior vaso que passa pelo ponto. A espessura desenhada é $r \times 10^{-3} \times$ exagero, convertendo milímetros em metros, com exagero ajustável (padrão $2\times$). Os vasos são desenhados do mais fino para o mais grosso, para que os troncos fiquem por cima nas bifurcações. A iluminação dos tubos (luz direcional e ambiente do pipeline fixo) é da API e serve apenas para dar volume; o estudo de iluminação é tema do TP2.

\textbf{Transformações e observador.} A matriz do modelo é $M = T \cdot R_x \cdot R_y \cdot R_z \cdot S \cdot N$, em que $N$ centraliza o modelo e normaliza seu tamanho para que a mesma câmera sirva a 2D e 3D. Há quatro vistas (frontal, lateral, superior e isométrica). Para que o modelo não mude de tamanho ao trocar de projeção, a caixa ortográfica tem meia altura $h = d \tan(\text{fovy}/2)$, igual à altura do tronco de pirâmide perspectivo na distância $d$ do alvo. Os eixos do mundo são desenhados só com $V$, permanecendo fixos enquanto o modelo é transformado.

\textbf{Antialiasing.} O SSAA do modo didático reutiliza o Bresenham e o Scanline: a cena é rasterizada numa grade $k$ vezes maior, com cada pixel de linha carimbado como um bloco $k \times k$ (para manter a largura de um pixel final), e reduzida com \texttt{reshape} e média. A mesma função com $k = 1$ gera a versão serrilhada, garantindo uma comparação justa. O MSAA do visualizador usa um framebuffer fora da tela (FBO) com 4 amostras por pixel, resolvido para a janela com \texttt{glBlitFramebuffer}, o que não depende de a janela oferecer multiamostragem.

\textbf{Recorte.} Os segmentos são recortados em coordenadas reais e só depois rasterizados, o que caracteriza recorte geométrico e não apenas descarte de pixels. O Sutherland--Hodgman aceita qualquer janela convexa, com o teste ``dentro'' feito pelo sinal do produto vetorial $(B - A) \times (P - A)$ para a janela em sentido anti-horário. No Weiler--Atherton, vértices exatamente sobre arestas, comuns em uma grade inteira, impedem a classificação de entrada e saída; nesses casos a janela é deslocada por $10^{-6}$ vezes o tamanho da cena, valor desprezível para o resultado.

% ======================================================================
\section{Resultados}

\subsection{Leitura e representação dos modelos}
A Tabela~\ref{tab:modelos} resume os modelos lidos. Em todos, o número de pontos é $2 N_{\text{term}}$ e o de segmentos é igual ao de pontos menos um, assinatura de uma árvore conexa e sem ciclos, o que confirma a leitura. A Figura~\ref{fig:m1} mostra as representações do visualizador.

\begin{table}[H]
  \centering
  \small
  \caption{Modelos CCO lidos: número de pontos, segmentos e intervalo do raio.}
  \label{tab:modelos}
  \begin{tabular}{lcccc|lcccc}
    \toprule
    \textbf{Modelo 2D} & \textbf{Pontos} & \textbf{Seg.} & $r_{\min}$ & $r_{\max}$ &
    \textbf{Modelo 3D} & \textbf{Pontos} & \textbf{Seg.} & $r_{\min}$ & $r_{\max}$ \\
    \midrule
    Nterm008 & 16  & 15  & 0,2491 & 0,6668 & Nterm016 & 32   & 31   & 0,2261 & 0,6870 \\
    Nterm016 & 32  & 31  & 0,2279 & 0,8176 & Nterm032 & 64   & 63   & 0,2132 & 0,8373 \\
    Nterm032 & 64  & 63  & 0,2058 & 0,9804 & Nterm064 & 128  & 127  & 0,2052 & 1,0178 \\
    Nterm064 & 128 & 127 & 0,2032 & 1,1814 & Nterm128 & 256  & 255  & 0,1454 & 1,2258 \\
    Nterm128 & 256 & 255 & 0,1454 & 1,1930 & Nterm256 & 512  & 511  & 0,1192 & 1,2472 \\
    Nterm256 & 512 & 511 & 0,1074 & 1,1937 & Nterm512 & 1024 & 1023 & 0,0956 & 1,2666 \\
    \bottomrule
  \end{tabular}
\end{table}

\begin{figure}[H]
  \centering
  \painel{0.24}{m1_linhas_jet}{(a) linhas, jet}\hfill
  \painel{0.24}{m1_faixas_jet}{(b) faixas, jet}\hfill
  \painel{0.24}{m1_faixas_viridis}{(c) faixas, viridis}\hfill
  \painel{0.24}{m1_tubos_3d}{(d) tubos 3D}
  \caption{Representações do modelo: (a) linhas coloridas pelo raio; (b) e (c) espessura e cor proporcionais ao raio nos dois mapas; (d) tubos 3D iluminados (Nterm256).}
  \label{fig:m1}
\end{figure}

\subsection{Transformações, observador e projeções}
A Figura~\ref{fig:m2t} mostra translação, rotação e escala aplicadas ao modelo, com os eixos do mundo fixos. A Figura~\ref{fig:m2v} mostra as quatro vistas do observador e a comparação entre projeções. Na perspectiva, o tronco próximo à câmera aparece mais largo que os ramos distantes de mesmo raio; na ortográfica, a espessura aparente não depende da profundidade. As matrizes próprias foram comparadas elemento a elemento com \texttt{glOrtho}, \texttt{gluPerspective}, \texttt{gluLookAt} e \texttt{glRotatef}, coincidindo em todos os casos (tolerância $10^{-5}$).

\begin{figure}[H]
  \centering
  \painel{0.24}{m2_original}{(a) original}\hfill
  \painel{0.24}{m2_translacao}{(b) translação}\hfill
  \painel{0.24}{m2_rotacao}{(c) rotação de $35^\circ$ em $z$}\hfill
  \painel{0.24}{m2_escala}{(d) escala $0{,}6$}
  \caption{Transformações geométricas aplicadas ao modelo. Os eixos do mundo (x em vermelho, y em verde) usam apenas a matriz de visão e permanecem fixos.}
  \label{fig:m2t}
\end{figure}

\begin{figure}[H]
  \centering
  \painel{0.155}{m2_vista_frontal}{(a) frontal}\hfill
  \painel{0.155}{m2_vista_lateral}{(b) lateral}\hfill
  \painel{0.155}{m2_vista_superior}{(c) superior}\hfill
  \painel{0.155}{m2_vista_isometrica}{(d) isométrica}\hfill
  \painel{0.155}{m2_perspectiva}{(e) perspectiva}\hfill
  \painel{0.155}{m2_ortografica}{(f) ortográfica}
  \caption{(a) a (d): quatro posições do observador para o modelo 3D Nterm256. (e) e (f): mesma câmera com projeção perspectiva e ortográfica (Nterm064).}
  \label{fig:m2v}
\end{figure}

\subsection{Rasterização de linhas e polígonos}
A Tabela~\ref{tab:bresenham} traz o exemplo numérico do Bresenham para o segmento $(20,10)$--$(30,18)$, com $\Delta x = 10$ e $\Delta y = 8$. Nos passos com $p_k < 0$ (passos 2 e 7), $y$ é mantido no passo seguinte. Em 20\,000 segmentos aleatórios, a implementação sempre começou e terminou nas extremidades, gerou pixels 8-conexos e manteve todos a no máximo meio pixel da reta ideal; o DDA divergiu em cerca de 17\% dos casos, sempre em empates de arredondamento, que são igualmente válidos.

\begin{table}[H]
  \centering
  \small
  \caption{Exemplo numérico do Bresenham de $(20,10)$ a $(30,18)$: $p_0 = 6$, $2\Delta y = 16$, $2\Delta y - 2\Delta x = -4$.}
  \label{tab:bresenham}
  \begin{tabular}{c|ccccccccccc}
    \toprule
    $k$ & 0 & 1 & 2 & 3 & 4 & 5 & 6 & 7 & 8 & 9 & 10 \\
    $p_k$ & 6 & 2 & $-2$ & 14 & 10 & 6 & 2 & $-2$ & 14 & 10 & 6 \\
    $x$ & 20 & 21 & 22 & 23 & 24 & 25 & 26 & 27 & 28 & 29 & 30 \\
    $y$ & 10 & 11 & 12 & 12 & 13 & 14 & 15 & 16 & 16 & 17 & 18 \\
    \bottomrule
  \end{tabular}
\end{table}

Para o Scanline Fill, o polígono côncavo $(2,1), (12,3), (9,7), (12,11), (5,9), (2,11)$ gera, na linha $y = 10$, quatro interseções ($2{,}00$; $3{,}50$; $8{,}50$; $11{,}25$) e dois trechos, $[2..3]$ e $[9..11]$, correspondentes à reentrância. Comparado a um teste par-ímpar ponto a ponto em 400 polígonos aleatórios, não houve divergências. A Figura~\ref{fig:raster} mostra a árvore rasterizada, o exemplo numérico e o preenchimento de uma região de interesse.

\begin{figure}[H]
  \centering
  \painel{0.24}{m3_arvore_bresenham}{(a) árvore, Bresenham}\hfill
  \painel{0.24}{m3_segmento_bresenham}{(b) segmento da Tab.~\ref{tab:bresenham}}\hfill
  \painel{0.24}{m4_scanline_animacao}{(c) Scanline em andamento}\hfill
  \painel{0.24}{m4_scanline}{(d) Scanline completo}
  \caption{Modo didático: (a) segmentos da árvore Nterm016 com a reta ideal sobreposta; (b) pixels do exemplo numérico; (c) e (d) região de interesse côncava preenchida por Scanline, com a linha de varredura em vermelho.}
  \label{fig:raster}
\end{figure}

\subsection{Aliasing e antialiasing}
A Figura~\ref{fig:aa} compara a mesma cena sem antialiasing e com SSAA $4\times$ (16 amostras por pixel), e o efeito do MSAA $4\times$ na borda de um tubo. No leque de linhas, a região central, em que as linhas ficam mais próximas que um pixel, vira um padrão quadriculado sem antialiasing e uma transição suave com SSAA.

\begin{figure}[H]
  \centering
  \painel{0.3}{m5_ssaa_leque}{(a) leque: sem AA $\times$ SSAA $4\times$}\hfill
  \painel{0.3}{m5_ssaa_arvore}{(b) árvore: sem AA $\times$ SSAA $4\times$}\hfill
  \painel{0.18}{m5_msaa_off_zoom}{(c) sem MSAA}\hfill
  \painel{0.18}{m5_msaa_on_zoom}{(d) MSAA $4\times$}
  \caption{(a) e (b): supersampling próprio, com a metade esquerda serrilhada e a direita suavizada. (c) e (d): borda de um tubo no visualizador, ampliada $3\times$.}
  \label{fig:aa}
\end{figure}

\subsection{Recorte de segmentos}
A Tabela~\ref{tab:clip} mostra o recorte do segmento $(2,6)$--$(36,30)$ pela janela $[10,30] \times [10,25]$. O Cohen--Sutherland calcula a interseção $(7{,}67;\ 10)$, que ainda está fora da janela (código \texttt{0001}), e precisa de uma nova interseção. O Liang--Barsky calcula quatro razões $q/p$ e apenas os dois pontos finais. Ambos chegam a $(10;\ 11{,}65)$--$(28{,}92;\ 25)$. Em 100\,000 segmentos aleatórios em $[-3,3]^2$ com janela $[-1,1]^2$, os resultados foram idênticos (50\,810 recortados, 47\,910 rejeitados e 1\,280 aceitos inteiros). A Figura~\ref{fig:clip} mostra o recorte da árvore.

\begin{table}[H]
  \centering
  \small
  \caption{Recorte do segmento $(2,6)$--$(36,30)$ pela janela $[10,30] \times [10,25]$.}
  \label{tab:clip}
  \begin{tabular}{lll|lrrrr}
    \toprule
    \multicolumn{3}{c|}{\textbf{Cohen--Sutherland}} & \multicolumn{5}{c}{\textbf{Liang--Barsky}} \\
    \textbf{Passo} & $P_0$ [código] & $P_1$ [código] & \textbf{Borda} & $p$ & $q$ & $t_0$ & $t_1$ \\
    \midrule
    0 & $(2{,}00;\ 6{,}00)$ [0101] & $(36{,}00;\ 30{,}00)$ [1010] & esquerda & $-34$ & $-8$ & 0,2353 & 1,0000 \\
    1 & $(7{,}67;\ 10{,}00)$ [0001] & $(36{,}00;\ 30{,}00)$ [1010] & direita & 34 & 28 & 0,2353 & 0,8235 \\
    2 & $(10{,}00;\ 11{,}65)$ [0000] & $(36{,}00;\ 30{,}00)$ [1010] & baixo & $-24$ & $-4$ & 0,2353 & 0,8235 \\
    3 & $(10{,}00;\ 11{,}65)$ [0000] & $(28{,}92;\ 25{,}00)$ [0000] & cima & 24 & 19 & 0,2353 & 0,7917 \\
    \bottomrule
  \end{tabular}
\end{table}

\begin{figure}[H]
  \centering
  \painel{0.32}{m6_recorte_antes}{(a) antes}\hfill
  \painel{0.32}{m6_recorte_depois}{(b) depois}\hfill
  \painel{0.32}{m6_recorte_sobreposto}{(c) sobreposto}
  \caption{Recorte dos segmentos da árvore Nterm032 pela janela em vermelho: 24 segmentos aceitos inteiros, 13 recortados e 26 rejeitados.}
  \label{fig:clip}
\end{figure}

\subsection{Recorte de polígonos}
No exemplo do Sutherland--Hodgman com o pentágono $(5,15), (20,3), (35,18), (22,32), (18,20)$ e a janela $[10,30] \times [10,25]$, o polígono passa de 5 para 6, 7, 8 e 9 vértices após cada aresta da janela. A Figura~\ref{fig:poly} mostra os dois casos que diferenciam os algoritmos. Na forma em U, cuja base fica fora da janela, o Sutherland--Hodgman devolve um polígono com uma ponte degenerada (em magenta), enquanto o Weiler--Atherton devolve dois polígonos (cores distintas). Com uma janela côncava em L, o Sutherland--Hodgman recorta pela interseção dos semiplanos das arestas e perde parte da região, enquanto o Weiler--Atherton produz o resultado correto. Em 600 pares de polígonos simples, côncavos e com casos degenerados, o Weiler--Atherton não apresentou divergências em 150\,000 pontos de teste.

\begin{figure}[H]
  \centering
  \painel{0.24}{m7_sh_u_depois}{(a) U, Sutherland--Hodgman}\hfill
  \painel{0.24}{m7_wa_u_depois}{(b) U, Weiler--Atherton}\hfill
  \painel{0.24}{m7_sh_l_concavo}{(c) janela em L, S--H}\hfill
  \painel{0.24}{m7_wa_l_concavo}{(d) janela em L, W--A}
  \caption{Recorte de polígonos. (a) ponte degenerada ligando as duas componentes; (b) duas componentes separadas; (c) resultado incorreto do Sutherland--Hodgman com janela côncava; (d) resultado correto do Weiler--Atherton.}
  \label{fig:poly}
\end{figure}

% ======================================================================
\section{Discussão}

\textbf{Comparação entre algoritmos de recorte de segmentos.} O Liang--Barsky evita interseções desperdiçadas e calcula apenas os pontos finais, enquanto o Cohen--Sutherland pode calcular interseções ainda externas, como no exemplo da Tabela~\ref{tab:clip}. Em compensação, a rejeição trivial do Cohen--Sutherland custa apenas operações sobre bits. No experimento com 100\,000 segmentos, em que quase metade foi rejeitada trivialmente, o Cohen--Sutherland foi mais rápido em Python (cerca de 0,7\,$\mu$s contra 1,1\,$\mu$s por segmento). A vantagem relativa depende, portanto, da proporção de segmentos totalmente fora da janela. O Cyrus--Beck seria necessário apenas para janelas convexas não retangulares.

\textbf{Sutherland--Hodgman e Weiler--Atherton.} O Sutherland--Hodgman é simples e eficiente, e o resultado preenchido por Scanline fica correto mesmo com pontes, porque elas têm área nula e ficam sobre linhas horizontais ou verticais. Porém o contorno exibe arestas espúrias e a saída não representa componentes separadas, o que é inadequado quando o resultado alimenta outras etapas. O Weiler--Atherton resolve esses casos e aceita janelas côncavas, ao custo de estruturas de dados mais complexas e de tratamento explícito de casos degenerados.

\textbf{Decisões de implementação.} A espessura foi modelada como geometria (faixas e tubos) e não com \texttt{glLineWidth}, que é limitado pelo hardware e aplica uma única largura por chamada. Os vértices são enviados uma única vez para a GPU em VBOs, pois o modelo Nterm512 gera 90\,088 triângulos. O MSAA foi feito por FBO porque o servidor gráfico do WSL não ofereceu janela multiamostrada.

\textbf{Limitações.} No ambiente de desenvolvimento (WSL com renderização por software), o modelo Nterm512 em tubos fica com taxa de quadros baixa durante a rotação. A iluminação dos tubos é da API. O deslocamento da janela no Weiler--Atherton resolve os casos degenerados encontrados, mas é uma solução numérica, não simbólica. O SSAA usa filtro caixa, que é o mais simples; filtros com pesos (como o gaussiano) suavizariam melhor.

% ======================================================================
\section{Conclusão}

O trabalho cumpriu os objetivos do TP1: os modelos CCO 2D e 3D foram lidos e representados por linhas, faixas e tubos codificados por cor e espessura; as transformações, as vistas e as projeções foram implementadas com matrizes próprias; e os algoritmos de rasterização, preenchimento, antialiasing e recorte foram implementados, validados e comparados, incluindo o Weiler--Atherton como diferencial. A separação entre o visualizador e o modo didático permitiu demonstrar cada algoritmo pixel a pixel sem depender da biblioteca. Como continuidade, o TP2 deve incorporar visibilidade, iluminação própria, transparência, textura e técnicas baseadas em raios sobre essa mesma base.

% --- REFERÊNCIAS BIBLIOGRÁFICAS ---
\begin{thebibliography}{99}
  \bibitem{schreiner1993} SCHREINER, W.; BUXBAUM, P. F. Computer-optimization of vascular trees. \textbf{IEEE Transactions on Biomedical Engineering}, v. 40, n. 5, p. 482--491, 1993.
  \bibitem{foley1990} FOLEY, J. D.; VAN DAM, A.; FEINER, S. K.; HUGHES, J. F. \textbf{Computer Graphics: Principles and Practice}. 2. ed. Reading: Addison-Wesley, 1990.
  \bibitem{hearn1997} HEARN, D.; BAKER, M. P. \textbf{Computer Graphics, C Version}. 2. ed. Upper Saddle River: Prentice Hall, 1997.
  \bibitem{bresenham1965} BRESENHAM, J. E. Algorithm for computer control of a digital plotter. \textbf{IBM Systems Journal}, v. 4, n. 1, p. 25--30, 1965.
  \bibitem{newman1979} NEWMAN, W. M.; SPROULL, R. F. \textbf{Principles of Interactive Computer Graphics}. 2. ed. New York: McGraw-Hill, 1979.
  \bibitem{liang1984} LIANG, Y.-D.; BARSKY, B. A. A new concept and method for line clipping. \textbf{ACM Transactions on Graphics}, v. 3, n. 1, p. 1--22, 1984.
  \bibitem{cyrus1978} CYRUS, M.; BECK, J. Generalized two- and three-dimensional clipping. \textbf{Computers \& Graphics}, v. 3, n. 1, p. 23--28, 1978.
  \bibitem{sutherland1974} SUTHERLAND, I. E.; HODGMAN, G. W. Reentrant polygon clipping. \textbf{Communications of the ACM}, v. 17, n. 1, p. 32--42, 1974.
  \bibitem{weiler1977} WEILER, K.; ATHERTON, P. Hidden surface removal using polygon area sorting. \textbf{ACM SIGGRAPH Computer Graphics}, v. 11, n. 2, p. 214--222, 1977.
  \bibitem{borland2007} BORLAND, D.; TAYLOR II, R. M. Rainbow color map (still) considered harmful. \textbf{IEEE Computer Graphics and Applications}, v. 27, n. 2, p. 14--17, 2007.
\end{thebibliography}

\end{document}
